\section{Versions as Interpretive Evidence}

A version is not merely a transparent window onto a single unchanged text.
Its wording can intensify, preserve, or soften the sequence on which an
interpretation depends. Four loci must be kept separate: the command to go or
look outside; the sun's approach to setting; completed sunset and darkness;
and the covenant's final temporal formula.

The Hebrew constellation of all four makes the daylight question sharp. The
Clementine Latin substantially preserves it. Swete's printed Greek keeps
enough of the movement to raise the issue, but its close is less emphatically
temporal. Augustine's Latin citation retains ``that day'' at verse 18. It
neither identifies his exemplar nor proves Greek variation, but it prevents
his Latin wording from being silently assimilated to Swete's spatial reading.

The Douay--Rheims is useful as an English witness to the Vulgate line, not as
an independent ancient authority. Modern readers can easily hear its sequence
as continuous. That readerly fact explains the present question; it does not
settle the textual history behind it.

One dated Syriac web witness preserves the outside, look, and count commands,
places the sun toward setting at verses 12 and 17, and retains ``on that day''
at verse 18 without an explicit interval. Because the site's textual base is
unidentified, this is evidence about that delivery, not a critical claim about
the Peshitta tradition. Jubilees 14 is not a version: as rewritten Scripture it
separates the star promise from the later sacrifice and reports the latter
sunset, but does not identify the hour of the star sight.

